\section{Conclusion}

This project provides a small and inspectable two-dimensional SPH solver that
follows the state-equation scheme of the course notes, with a stiff state
equation, mirrored boundary pressures, adaptive time stepping, an interactive
visualisation, and a reproducible experiment pipeline.

With the stiffness coefficient chosen so that pressure waves are ten times
faster than the free-fall velocity, the settled hydrostatic column keeps its
density within about $0.7\,\%$ of the rest density and reproduces the
hydrostatic pressure to within a few percent when averaged over time. The dam
break stays below $3\,\%$ density error until the water hits the opposite wall.

The parameter studies show that the stiffness coefficient trades
compressibility against cost and pressure noise, and that the smoothing length
should stay close to $1.3$ particle spacings, because larger values trigger the
pairing instability. Mirroring the pressure to the boundary samples removes the
gap between the wall and the first fluid layers, but it feeds energy into the
fluid: a viscosity of at least $\nu = 600$ is needed to keep the fluid stable,
which slows the flow, and slow eddies remain along the walls.

Future work could extrapolate the boundary pressure instead of mirroring it,
replace the all-pairs neighbour search by a uniform grid to reach higher
resolutions, reduce the pressure noise with an iterative or implicit pressure
solver (PCISPH, IISPH, DFSPH), and compare the dam-break front quantitatively
with the measurements of Martin and Moyce.
